\input{../tex-inputs/latex-korrekturansicht-vorspann}

\section[Arthur Schnitzler an Gustav Schwarzkopf, 1. 7. 1925]{L04361 Arthur Schnitzler an Gustav Schwarzkopf, 1. 7. 1925}
\nopagebreak\mylabel{L04361v}
\rehead{ }\normalsize\beginnumbering\briefempfaengerindex{Schwarzkopf, Gustav@\textsc{Schwarzkopf, Gustav}!zzzSchnitzler, Arthur@\emph{von Arthur Schnitzler}!1925-07-011@{1. 7. 1925}|(be}
\toendnotes[C]{\smallbreak\pagebreak[2]}
\correspDesc{Versand  durch Arthur Schnitzler am 1. 7. 1925 in Bozen
\newline{}Übermittlung  am 2. 7. 1925 in Bozen
\newline{}Erhalt  durch Gustav Schwarzkopf im Zeitraum [2. 7. 1925 – 6. 7. 1925?] in Wien}\toendnotes[C]{\smallbreak}
\Standort{DLA, A:Schnitzler, HS.1985.1.1897.}
\physDesc{Bildpostkarte, 221 Zeichen
\newline{}Handschrift: Bleistift, lateinische Kurrent
\newline{}Versand: Stempel: »\nobreak{}\oindex{Bozen@\textbf{Bozen}, \emph{Hauptstadt}|pwk}Bolzano (Arrivi e partenze), 2–7. 25, 11\nobreak{}«.  }\toendnotes[C]{\smallbreak}\pstart{}{\pb}Hrn Gustav\pend{}\pstart{}Schwarzkopf\pend{}\pstart{}\textcolor{pink}{Wien I}\oindex{I., Innere Stadt@\textbf{I., Innere Stadt}, \emph{Verwaltungsgebiet}|pw}{}\ledrightnote{\textcolor{pink}{I., Innere Stadt}}\pend{}\pstart{}\textcolor{pink}{Tiefer Graben 17}\oindex{Wien@\textbf{Wien}!I., Innere Stadt@\textbf{I., Innere Stadt}!Tiefer Graben 17@\textbf{Tiefer Graben 17}, \emph{Wohngebäude}|pw}{}\ledrightnote{\textcolor{pink}{Tiefer Graben 17}}\pend{}{\bigskip}
\pstart
           \noindent{}\centering{}{\pb}\textcolor{gray}{\textbf{\textcolor{pink}{Bolzano: Piazza Walter}\oindex{Waltherplatz@\textbf{Waltherplatz}, \emph{Platz}|pw}{}\ledrightnote{\textcolor{pink}{Waltherplatz}}.}}\pend
           \vspace{1em}
\pstart
           \raggedleft{}{\pb}1. 7. 2\textcolor{gray}{5}.\pend
           \vspace{0.5em}
\pstart
           Herzliche Grüße, und
      auf baldgs Wiedersehen!
      (Daſs ich noch \uline{nicht} in
      \textcolor{pink}{Wien}\oindex{Wien@\textbf{Wien}, \emph{Verwaltungsgebiet}|pw}{}\ledrightnote{\textcolor{pink}{Wien}}, konnten Sie wol
               aus der \label{K_L04361-1v}\edtext{\textcolor{green}{Notiz}\pwindex{?? [Zeitungsnotiz, dass Schnitzler nicht zur Premiere von Schleier der Pierrette kommt]@\emph{?? [Zeitungsnotiz, dass Schnitzler nicht zur Premiere von Schleier der Pierrette kommt]}|pwv}{}\ledrightnote{{$\rightarrow$}\emph{\textcolor{green}{?? [Zeitungsnotiz, dass Schnitzler nicht zur Premiere von Schleier der Pierrette kommt]}}}}{\lemma{\textnormal{\emph{Notiz}}}\Cendnote{\textnormal{nicht nachgewiesen; entsprechend kryptisch bleibt auch die 
               diesbezügliche Formulierung \textcolor{blue}{Schnitzlers}.}}}\label{K_L04361-1} entnehmen,
               daſs ich der \label{K_L04361-2v}\edtext{\textcolor{violet}{Première}\eventindex{Volkstheater@\textbf{Volkstheater}!Premiere von Der Schleier der Pierrette, 3.7.1925@Premiere von Der Schleier der Pierrette, 3.7.1925|pwv}{}\ledrightnote{{$\rightarrow$}\emph{\textcolor{violet}{Premiere von Der Schleier der Pierrette, 3.7.1925}}}}{\lemma{\textnormal{\emph{Première}}}\Cendnote{\textnormal{
                  Die \textcolor{violet}{Premiere der Pantomime \emph{\textcolor{green}{Der Schleier der Pierrette}\pwindex{Schnitzler, Arthur 15. 5. 1862 Wien – 21. 10. 1931 ebd.@\textsc{Schnitzler, Arthur} (15. 5. 1862 Wien – 21. 10. 1931 ebd.), \emph{Schriftsteller, Mediziner}!Schleier der Pierrette. Pantomime in drei Bildern@\strich\emph{Der Schleier der Pierrette. Pantomime in drei Bildern}|pwk}} von \textcolor{blue}{Arthur Schnitzler}}\eventindex{Volkstheater@\textbf{Volkstheater}!Premiere von Der Schleier der Pierrette, 3.7.1925@Premiere von Der Schleier der Pierrette, 3.7.1925|pwk} fand am 3. 7. 1925 im \textcolor{pink}{Deutschen Volkstheater}\oindex{Wien@\textbf{Wien}!VII., Neubau@\textbf{VII., Neubau}!Volkstheater@\textbf{Volkstheater}, \emph{Theater}|pwk} statt, die \textcolor{violet}{2. Aufführung}\eventindex{Volkstheater@\textbf{Volkstheater}!2. Aufführung von Der Schleier der Pierrette, 5.7.1925@2. Aufführung von Der Schleier der Pierrette, 5.7.1925|pwkv} zwei Tage danach.
                  Es handelte sich um ein Gastspiel des \emph{\textcolor{brown}{Kamerny Theatre}\orgindex{Kamerny Theatre@Kamerny Theatre|pwk}} mit \textcolor{blue}{Alexander J. Tairow}\pwindex{Tairow, Alexander J. 24.\,6.\,1885 Romny – 27.\,9.\,1950 Moskau@\textsc{Tairow, Alexander J.} (24.\,6.\,1885 Romny – 27.\,9.\,1950 Moskau), \emph{Theaterleiter, Regisseur, Schauspieler}|pwk}. \textcolor{blue}{Schnitzler}
                  nahm nur an der zweiten (und letzten) Vorstellung teil, vgl. A. S.: \emph{Kulturveranstaltungen}, 5. 7. 1925.
               }}}\label{K_L04361-2} der
      \textcolor{green}{Pierette}\pwindex{Schnitzler, Arthur 15. 5. 1862 Wien – 21. 10. 1931 ebd.@\textsc{Schnitzler, Arthur} (15. 5. 1862 Wien – 21. 10. 1931 ebd.), \emph{Schriftsteller, Mediziner}!Schleier der Pierrette. Pantomime in drei Bildern@\strich\emph{Der Schleier der Pierrette. Pantomime in drei Bildern}|pw}{}\ledrightnote{\textcolor{green}{Der Schleier der Pierrette. Pantomime in drei Bildern}} beiwohnen würde.)\pend
           
\pstart
           Ihr{\\[\baselineskip]}\spacefill\mbox{Arthur}\pend
           \leftskip=0em{}\selectlanguage{ngerman}\endnumbering\briefempfaengerindex{Schwarzkopf, Gustav@\textsc{Schwarzkopf, Gustav}!zzzSchnitzler, Arthur@\emph{von Arthur Schnitzler}!1925-07-011@{1. 7. 1925}|)be}\mylabel{L04361h}
\begin{anhang}
\end{anhang}\newcommand{\dateiname}{L04361}\newcommand{\titel}{Arthur Schnitzler an Gustav Schwarzkopf, 1. 7. 1925}\newcommand{\editorInnen}{Herausgegeben von Jahnke, SelmaMüller, Martin Anton}\input{../tex-inputs/latex-korrekturansicht-abspann}
